\documentclass[11pt,a4paper]{article}

\usepackage[margin=1in]{geometry}
\usepackage[T1]{fontenc}
\usepackage[utf8]{inputenc}

\usepackage{microtype}
\usepackage{xcolor}
\usepackage{graphicx}
\usepackage{hyperref}
\usepackage{enumitem}
\usepackage{array}
\usepackage{tabularx}
\usepackage{booktabs}
\usepackage{longtable}
\usepackage{float}
\usepackage{ifthen}
\usepackage{titlesec}
\usepackage{fancyhdr}
\usepackage{tcolorbox}
\usepackage{amssymb}

\hypersetup{
  colorlinks=true,
  linkcolor=blue!50!black,
  urlcolor=blue!50!black,
  pdftitle={READCRAFTERY — Artist Brief (Post-M1 Technical Handoff)},
  pdfauthor={OpenAI},
  pdfsubject={Artist production handoff},
  pdfcreator={LaTeX}
}

\definecolor{RCBrown}{HTML}{2C1A0E}
\definecolor{RCParchment}{HTML}{F5E6C8}
\definecolor{RCGold}{HTML}{D4A017}
\definecolor{RCAmber}{HTML}{E8721A}
\definecolor{RCMagic}{HTML}{7B4FBE}
\definecolor{RCBlue}{HTML}{1E3A5F}
\definecolor{RCCorrect}{HTML}{4CAF50}
\definecolor{RCIncorrect}{HTML}{E53935}
\definecolor{RCHint}{HTML}{FF9800}
\definecolor{RCNeutral}{HTML}{9E9E9E}
\definecolor{RCSection}{HTML}{3E2A1F}
\definecolor{RCLightBorder}{HTML}{D8C6A3}
\definecolor{RCSoftBg}{HTML}{FBF7EF}

\pagestyle{fancy}
\fancyhf{}
\lhead{READCRAFTERY — Artist Brief}
\rhead{Post-M1 Technical Handoff}
\cfoot{\thepage}
\setlength{\headheight}{14pt}

\titleformat{\section}{\large\bfseries\color{RCSection}}{\thesection.}{0.5em}{}
\titleformat{\subsection}{\normalsize\bfseries\color{RCSection}}{\thesubsection}{0.5em}{}
\titleformat{\subsubsection}{\normalsize\bfseries}{\thesubsubsection}{0.5em}{}

\setlist[itemize]{topsep=4pt,itemsep=2pt,leftmargin=1.6em}
\setlist[enumerate]{topsep=4pt,itemsep=2pt,leftmargin=1.8em}

\newcommand{\docversion}{0.2}
\newcommand{\docdate}{March 2026}

% ---------- Image path placeholders ----------
% Replace these paths with your actual image files.
% If a file does not exist, the PDF will show a framed placeholder box instead.

% Canon Owlorumo reference — source of truth for silhouette and proportions.
\newcommand{\CanonCharacterPath}{path/to/canon_owlorumo_reference.png}

% Approved concept art — mood, texture, and emotional tone.
\newcommand{\ConceptArtPath}{path/to/approved_owlorumo_concept_art.png}

% PassageView mock — size and hierarchy reference.
\newcommand{\PassageViewMockPath}{path/to/passageview_context_mock.png}

% Library mock — scene-role reference.
\newcommand{\LibraryMockPath}{path/to/library_context_mock.png}

% Mini moodboard image — optional 2–3 image collage.
\newcommand{\MoodboardPath}{path/to/mini_moodboard.png}

% Do / Don't strip — optional alignment panel.
\newcommand{\DoDontPath}{path/to/do_dont_strip.png}

% ---------- Utility macros ----------

% Prints a color swatch and the hex code using the encoded color.
% Usage: \hexcode{7B4FBE}
\newcommand{\hexcode}[1]{%
  \begingroup
  \definecolor{temphex}{HTML}{#1}%
  \fcolorbox{black}{temphex}{\rule{0pt}{1.1ex}\rule{1.6ex}{0pt}}\hspace{0.45em}%
  {\ttfamily\textcolor{temphex}{\##1}}%
  \endgroup
}

\newcommand{\checkbox}{\ensuremath{\square}}

\newcommand{\briefbox}[2]{%
  \begin{tcolorbox}[
    colback=RCSoftBg,
    colframe=RCLightBorder,
    boxrule=0.5pt,
    arc=1mm,
    left=4mm,right=4mm,top=2.5mm,bottom=2.5mm,
    title=\textbf{#1},
    coltitle=black,
    fonttitle=\bfseries
  ]
  #2
  \end{tcolorbox}
}

\newcommand{\imageref}[4]{%
  \begin{figure}[H]
    \centering
    \IfFileExists{#1}{%
      \includegraphics[width=0.92\linewidth]{#1}%
    }{%
      \fbox{\begin{minipage}[c][52mm][c]{0.88\linewidth}
      \centering
      \textit{Image placeholder}\\[0.6em]
      {\ttfamily\detokenize{#1}}\\[0.8em]
      \small #4
      \end{minipage}}%
    }
    \caption{#2}
    \label{#3}
  \end{figure}
}

\begin{document}

\begin{titlepage}
\centering
\vspace*{2cm}

{\Huge\bfseries READCRAFTERY\\[0.4em]Artist Brief\par}
\vspace{0.3cm}
{\Large Post-M1 Technical Handoff\par}
\vspace{1cm}

\begin{tcolorbox}[
  width=0.9\linewidth,
  colback=RCParchment,
  colframe=RCLightBorder,
  boxrule=0.7pt,
  arc=1mm,
  auto outer arc
]
\centering
\textbf{Version:} \docversion \quad | \quad
\textbf{Date:} \docdate \quad | \quad
\textbf{Status:} Internal production brief\\[0.5em]
\textbf{Audience:} External artist / internal art production\\[0.5em]
\textbf{Canonical references:} Art Style Guide v1.5 · Asset Contracts (synced) · Architecture Contract v1.3 · Canon Core (synced) · Roadmap v5
\end{tcolorbox}

\vfill

{\large This file is intentionally production-oriented.\\
It assumes the technical slice already exists and the art layer is replacing placeholders, not redefining the system.\par}

\vfill
\end{titlepage}

\section{Purpose of this brief}

This is \textbf{not} the long-form style bible. This is the production handoff for the moment when \textbf{M1 technical is already complete and stable}.

At this stage:
\begin{itemize}
  \item all target screens already exist in-engine,
  \item all final node types are already in place,
  \item placeholder colors already match asset contracts,
  \item scene layout is already locked enough that real art can replace placeholders.
\end{itemize}

The artist is \textbf{not} asked to invent system behavior, layout, or UX. The artist is asked to deliver final production art that drops into a working technical slice.

\section{What this brief is for}

The goal of this handoff is to produce the \textbf{first real visual pass} for READCRAFTERY after the technical milestone is proven.

That means:
\begin{itemize}
  \item replace M0/M1 placeholders with real production-ready art,
  \item preserve all locked dimensions, anchors, and node types,
  \item keep readability and child clarity above decoration,
  \item produce only the assets required for the \textbf{M1 final vertical slice}.
\end{itemize}

This brief does \textbf{not} authorize broader worldbuilding, alternate themes, extra VFX, or speculative polish.

\section{Production context}

READCRAFTERY is a reading game for children aged \textbf{4–9}. The core mechanic is reading-first, not action-first.

The visual world is called \textbf{Storybook Arcane}:
\begin{itemize}
  \item warm,
  \item alive,
  \item handmade,
  \item magical but not explosive,
  \item child-safe and welcoming.
\end{itemize}

Two source inspirations define the feel:
\begin{itemize}
  \item \textbf{Choose Your Own Adventure} → each screen feels like a real book page.
  \item \textbf{Dr. Seuss} → organic shape language, expressive personality, invented texture.
\end{itemize}

Avoid completely:
\begin{itemize}
  \item modern flat SaaS look,
  \item anime/chibi proportions,
  \item realism,
  \item dark fantasy,
  \item polished vector smoothness.
\end{itemize}

\section{Critical production rule}

\briefbox{Non-negotiable rule}{
\textbf{This handoff begins only after M1 technical is complete.}

Before that moment:
\begin{itemize}
  \item M0 and technical M1 use placeholders,
  \item concept art is reference only,
  \item no final atlas production is expected,
  \item no art should block architecture, loading, save flow, or HTML5 validation.
\end{itemize}

The artist must assume the engine side already works. The art layer must be a \textbf{drop-in replacement}, not a redesign request.
}

\section{Single source of truth hierarchy}

If documents ever conflict, this order wins:
\begin{enumerate}
  \item \textbf{Asset Contracts} — dimensions, node types, anchors, states, placeholder replacement logic.
  \item \textbf{Art Style Guide} — visual style, palette, character identity, rendering rules.
  \item \textbf{Architecture Contract} — runtime rules and invariants.
  \item \textbf{Roadmap} — production sequencing and milestone timing.
  \item This brief — applied summary for art production.
\end{enumerate}

If there is still doubt, stop and ask before producing.

\section{Core visual direction}

\subsection{Style name}
\textbf{Storybook Arcane}

\subsection{Emotional target}
The game should feel like:
\begin{itemize}
  \item a magical library on a rainy evening,
  \item a safe storybook world,
  \item handcrafted and warm,
  \item visually rich, but never at the cost of reading clarity.
\end{itemize}

\subsection{Most important rule}
\textbf{Text is king.} If an art decision makes the text less readable, that art decision is wrong.

\section{Owlorumo — character production summary}

Owlorumo is the official companion character. He is a \textbf{young owl wizard}: kind, curious, encouraging, never intimidating.

\subsection{Must-read silhouette}
Owlorumo must read clearly as:
\begin{itemize}
  \item owl,
  \item wizard hat,
  \item staff with glowing crystal.
\end{itemize}

\subsection{Key identity features}
\begin{itemize}
  \item large amber eyes,
  \item small curved beak,
  \item navy robe,
  \item wrapped neck cloth,
  \item crescent moon clasp,
  \item bent wizard hat,
  \item brown boots,
  \item dark twisted staff,
  \item purple-white crystal glow (idle tone: \hexcode{7B4FBE}).
\end{itemize}

\subsection{Hard rules}
\begin{itemize}
  \item no beard,
  \item never threatening,
  \item never smug or condescending,
  \item staff is integral, not optional,
  \item silhouette must remain legible at native sprite size.
\end{itemize}

\section{Pixel art production rules}

\subsection{Rendering rules}
\begin{itemize}
  \item true pixel art,
  \item no smooth digital painting,
  \item no vector-clean curves,
  \item no blur,
  \item no gradient shading,
  \item no manual anti-aliasing,
  \item no pillow shading.
\end{itemize}

\subsection{Outline rules}
\begin{itemize}
  \item 1px outline only,
  \item never pure black,
  \item use darker local color instead.
\end{itemize}

\subsection{Lighting rules}
\begin{itemize}
  \item light direction: upper-left,
  \item flat shadow blocks only,
  \item max 2 shadow values per element,
  \item highlight used sparingly.
\end{itemize}

\subsection{Palette examples}
The final production palette is defined in the Art Style Guide. In this PDF, any hex code can be shown using its own encoded color. Examples:
\begin{itemize}
  \item Background dark: \hexcode{2C1A0E}
  \item Parchment: \hexcode{F5E6C8}
  \item Reward gold: \hexcode{D4A017}
  \item Magic glow: \hexcode{7B4FBE}
  \item Correct feedback: \hexcode{4CAF50}
  \item Error feedback: \hexcode{E53935}
\end{itemize}

\section{Technical size rules}

These are locked.

\begin{longtable}{@{}p{0.36\linewidth}p{0.22\linewidth}p{0.22\linewidth}p{0.14\linewidth}@{}}
\toprule
\textbf{Asset} & \textbf{Native size} & \textbf{Display} & \textbf{Scale} \\
\midrule
Owlorumo (PassageView) & 64×64 px & 192×192 px & 3× \\
Owlorumo (Library/menu) & 64×64 px & 256×256 px & 4× \\
Owlorumo (Celebration) & 64×64 px & 256–320 px & 4–5× \\
Accessories (overlay) & 64×64 px & matches context & same as base \\
UI icons & 16×16 px & 64×64 px & 4× \\
Buttons / tiles & 16×16 px base & 64×64 px & 4× \\
Buttons (touch-driven where specified) & 24×24 px min & 96×96 px target & 4× \\
Background layers & 480×270 px & 1920×1080 px & 4× \\
Parchment / wood tiles & 64×64 px & repeated in UI & 4× \\
Stars small & 8×8 px & 32×32 px & 4× \\
Stars large & 16×16 px & 64×64 px & 4× \\
\bottomrule
\end{longtable}

\textbf{Do not upscale before export.} Deliver at native draw size only.

\section{Engine-facing rules}

The engine is already built against the contract. The artist must \textbf{not} change:
\begin{itemize}
  \item node types,
  \item layout role,
  \item anchor point assumptions,
  \item state names,
  \item naming convention,
  \item expected frame counts unless approved.
\end{itemize}

Art replaces textures. It does not redefine UI architecture.

\section{Delivery format}

\subsection{Required source delivery}
\begin{itemize}
  \item preferred: \texttt{.aseprite}
  \item acceptable: layered \texttt{.psd}
\end{itemize}

\subsection{Required export delivery}
\begin{itemize}
  \item PNG with transparency where applicable,
  \item one file per state or frame as requested,
  \item no embedded scale-up,
  \item no texture atlas packing unless explicitly requested later.
\end{itemize}

\subsection{Naming}
Use the contract names exactly. Examples:
\begin{itemize}
  \item \texttt{owl\_idle\_01.png}
  \item \texttt{owl\_thinking\_03.png}
  \item \texttt{tile\_word\_found.png}
  \item \texttt{bg\_library\_layer\_4.png}
\end{itemize}

\subsection{Registration / anchor}
\begin{itemize}
  \item Owlorumo anchor: center-bottom,
  \item accessories: pixel-perfect registration with base sprite,
  \item UI assets: preserve exact margins for NineSlice where applicable.
\end{itemize}

\section{Approval sequence}

To minimize rework, approval happens in this order:
\begin{enumerate}
  \item Owlorumo silhouette sheet,
  \item Owlorumo color pass,
  \item Owlorumo idle,
  \item Owlorumo thinking,
  \item Owlorumo hint,
  \item PassageView textures and WordGlow UI tiles,
  \item Library background layers,
  \item Remaining M1-final states.
\end{enumerate}

Do not fully produce later items before earlier ones are approved.

\section{What the artist is producing in this handoff}

This handoff is for the \textbf{M1 final vertical slice}. It is intentionally limited.

\subsection{In scope}
\begin{itemize}
  \item Owlorumo production sprite work for approved states,
  \item PassageView textures,
  \item WordGlow UI tiles and icons,
  \item minimum library background art required for the slice,
  \item celebration stub replacement only if the milestone calls for it.
\end{itemize}

\subsection{Not in scope}
\begin{itemize}
  \item alternate themes,
  \item accessories system,
  \item community mod art,
  \item RTL assets,
  \item speculative polish,
  \item VFX-heavy celebration sequences,
  \item book perspective system beyond approved slice needs,
  \item final atlas optimization pass.
\end{itemize}

\section{Visual quality targets}

The final art should satisfy all of these:
\begin{itemize}
  \item child-safe,
  \item readable at gameplay size,
  \item warm, not sterile,
  \item magical, not explosive,
  \item expressive but not noisy,
  \item screen hierarchy always favors text and puzzle interaction.
\end{itemize}

If a beautiful asset competes with the reading task, it must be simplified.

\section{Acceptance checks for artist review}

Before approving an asset, verify:
\begin{itemize}
  \item it matches locked draw size,
  \item silhouette reads clearly at gameplay size,
  \item outline is consistently 1px,
  \item light direction is upper-left,
  \item no blur / no anti-aliasing / no gradients,
  \item palette remains in-family with project tokens,
  \item readability remains stronger than decoration,
  \item anchor assumptions still hold,
  \item export name matches contract exactly.
\end{itemize}

\section{Visual Reference Pack}

The brief should include a \textbf{small, controlled visual pack}. Do not send a giant moodboard. Do not send conflicting references. Every included image must have a clear role.

\briefbox{Image handling rules}{
\begin{itemize}
  \item Keep the pack small: \textbf{4 to 6 images total}.
  \item Do not include deprecated concept art.
  \item Do not include contradictory old versions.
  \item Do not include placeholder screenshots unless clearly labeled as layout-only.
  \item Label every image with its role.
  \item If a conflict appears, the \textbf{canon character reference wins}.
\end{itemize}
}

% PLACEHOLDER: Canon character reference — source of truth for Owlorumo.
\imageref{\CanonCharacterPath}{\textbf{Canon character reference — source of truth for Owlorumo.} If any other reference conflicts with this one, this image wins.}{fig:canon}{Use this slot for the single most authoritative Owlorumo image.}

% PLACEHOLDER: Approved concept art — mood, texture, and emotional tone.
\imageref{\ConceptArtPath}{\textbf{Approved concept art — use for mood, texture, and emotional tone.} Do not override the canon character reference on silhouette or key proportions.}{fig:concept}{Use this slot for the main approved concept art.}

% PLACEHOLDER: PassageView context mock — size and hierarchy reference.
\imageref{\PassageViewMockPath}{\textbf{PassageView context mock — size and hierarchy reference.} Text and puzzle readability always outrank character decoration.}{fig:passageview}{Use this slot for a layout mock or high-level composition mock of PassageView.}

% PLACEHOLDER: Library context mock — scene-role reference.
\imageref{\LibraryMockPath}{\textbf{Library context mock — scene-role reference.} Available books must stand out immediately; Owlorumo supports the scene but does not dominate selection.}{fig:library}{Use this slot for a library composition mock.}

% PLACEHOLDER: Mini moodboard — 2 to 3 images maximum, ideally combined into one collage.
\imageref{\MoodboardPath}{\textbf{Moodboard — atmosphere only.} Use for emotional direction, not for overriding the project’s locked character design or technical pixel-art rules.}{fig:moodboard}{Use this slot for a compact 2–3 image collage only.}

% PLACEHOLDER: Do / Don't strip — optional but recommended.
\imageref{\DoDontPath}{\textbf{Do / Don't strip — fast alignment tool.} If a new sketch drifts toward any “Don't” column, correct direction before production continues.}{fig:dodont}{Use this slot for a small yes/no comparison strip.}

\subsection{Recommended order inside the brief}
Put the images in this order:
\begin{enumerate}
  \item canon character reference,
  \item approved concept art,
  \item PassageView context mock,
  \item Library context mock,
  \item mini moodboard,
  \item Do / Don't strip.
\end{enumerate}

This order matters. It teaches the artist:
\begin{itemize}
  \item who the character is,
  \item how the character should feel,
  \item where the character lives in UI,
  \item what atmosphere to target,
  \item what styles to avoid.
\end{itemize}

\section{Final note to artist}

READCRAFTERY is not asking for maximal visual density. It is asking for \textbf{clarity with warmth}. The best result is not the most elaborate result. The best result is the one that makes a child want to read the next word.

\newpage
\section*{READCRAFTERY — Asset Checklist (Post-M1 Technical Handoff)}
\addcontentsline{toc}{section}{READCRAFTERY — Asset Checklist (Post-M1 Technical Handoff)}

\subsection*{A. Handoff assumptions}
Before production starts, all of these must already be true:
\begin{itemize}
  \item[\checkbox] M1 technical milestone complete.
  \item[\checkbox] HTML5 slice runs correctly.
  \item[\checkbox] Placeholders already wired to final node types.
  \item[\checkbox] Owlorumo size contract frozen at 64×64 master draw size.
  \item[\checkbox] Asset Contracts accepted as dimension/source-of-truth document.
  \item[\checkbox] Scene layout stable enough that art replacement will not cause redesign.
\end{itemize}

If any box above is unchecked, do not start external asset production.

\subsection*{B. Owlorumo — character assets}
\textbf{Required first approvals}
\begin{itemize}
  \item[\checkbox] Owlorumo silhouette sheet
  \item[\checkbox] Owlorumo color pass
  \item[\checkbox] Owlorumo expression sheet
\end{itemize}

\textbf{Required state deliveries for first production pass}
\begin{itemize}
  \item[\checkbox] \texttt{icon\_hint.png} — staff crystal motif, 16×16px. UI indicator only.
  Note: Owlorumo himself is the hint button. There is no separate btn\_hint asset.
  \item[\checkbox] \texttt{owl\_idle\_02.png}
  \item[\checkbox] \texttt{owl\_idle\_03.png}
  \item[\checkbox] \texttt{owl\_idle\_04.png}
  \item[\checkbox] \texttt{owl\_thinking\_01.png}
  \item[\checkbox] \texttt{owl\_thinking\_02.png}
  \item[\checkbox] \texttt{owl\_thinking\_03.png}
  \item[\checkbox] \texttt{owl\_thinking\_04.png}
  \item[\checkbox] \texttt{owl\_hint\_01.png}
  \item[\checkbox] \texttt{owl\_hint\_02.png}
  \item[\checkbox] \texttt{owl\_hint\_03.png}
  \item[\checkbox] \texttt{owl\_hint\_04.png}
\end{itemize}

\textbf{Optional later in same handoff, only if approved}
\begin{itemize}
  \item[\checkbox] \texttt{owl\_happy\_01.png} to \texttt{owl\_happy\_06.png}
  \item[\checkbox] \texttt{owl\_reading\_01.png} to \texttt{owl\_reading\_04.png}
  \item[\checkbox] \texttt{owl\_wave\_01.png} to \texttt{owl\_wave\_06.png}
\end{itemize}

\textbf{Deferred}
\begin{itemize}
  \item[\checkbox] excited
  \item[\checkbox] celebrate
  \item[\checkbox] sleepy
  \item[\checkbox] accessories
  \item[\checkbox] state\_4 fulfilled restoration (blue mantle illuminated white-gold, same identity anchors)
\end{itemize}

\subsection*{C. PassageView textures}
\begin{itemize}
  \item[\checkbox] \texttt{parchment\_tile.png} — 64×64 tileable
  \item[\checkbox] \texttt{wood\_tile.png} — 64×64 tileable
\end{itemize}

\textbf{Checks}
\begin{itemize}
  \item[\checkbox] parchment contrast remains low-noise behind text
  \item[\checkbox] wood texture does not compete with puzzle readability
  \item[\checkbox] both tile seamlessly
\end{itemize}

\subsection*{D. WordGlow UI assets}
\begin{itemize}
  \item[\checkbox] \texttt{tile\_word\_normal.png}
  \item[\checkbox] \texttt{tile\_word\_selected.png}
  \item[\checkbox] \texttt{tile\_word\_found.png}
  \item[\checkbox] \texttt{slot\_word\_empty.png} if included in approved slice
  \item[\checkbox] \texttt{slot\_word\_filled.png} if included in approved slice
\end{itemize}

\textbf{Checks}
\begin{itemize}
  \item[\checkbox] NineSlice corners survive scaling
  \item[\checkbox] tile state differences are readable instantly
  \item[\checkbox] found state feels rewarded, not noisy
\end{itemize}

\subsection*{E. Global UI icons / buttons}
\begin{itemize}
  \item[\checkbox] \texttt{icon\_hint.png} — staff crystal motif, 16×16px.
    No \texttt{btn\_hint.png} — Owlorumo himself is the hint button.
  \item[\checkbox] \texttt{btn\_exit.png} — door motif, 24×24px. Tap-hold 1 second to activate.
  \item[\checkbox] \texttt{book\_tts.png} — small open book, 16×16px. TTS play button in PassageView.
  \item[\checkbox] \texttt{book\_settings.png} — closed book, navy \#1E2A4A with crescent moon on spine, 16×24px. Fixed top-right of Library.
\end{itemize}

\textbf{Optional if included in current slice}
\begin{itemize}
  \item[\checkbox] \texttt{btn\_continue.png}
  \item[\checkbox] \texttt{btn\_back.png}
\end{itemize}

\textbf{Checks}
\begin{itemize}
  \item[\checkbox] icon silhouettes read at 16×16 native
  \item[\checkbox] button artwork respects touch target rules in engine
  \item[\checkbox] exit door reads as "way out" without text for a 5-year-old
  \item[\checkbox] settings book is visually distinct from regular library books
\end{itemize}

\subsection*{F. Library slice assets}
\textbf{Minimum background delivery}
\begin{itemize}
  \item[\checkbox] \texttt{bg\_library\_layer\_4.png}
  \item[\checkbox] \texttt{bg\_library\_layer\_3.png}
\end{itemize}

\textbf{Optional if current slice needs them}
\begin{itemize}
  \item[\checkbox] \texttt{bg\_library\_layer\_2.png}
  \item[\checkbox] \texttt{bg\_library\_layer\_1.png}
  \item[\checkbox] \texttt{book\_available.png}
  \item[\checkbox] \texttt{book\_dormant.png}
  \item[\checkbox] \texttt{book\_lock\_icon.png}
\end{itemize}

\textbf{Checks}
\begin{itemize}
  \item[\checkbox] available books stand out immediately
  \item[\checkbox] dormant vs readable is understandable in under one second
  \item[\checkbox] Owlorumo does not visually dominate book selection
\end{itemize}

\subsection*{G. Celebration / transition assets}
\textbf{Minimum only}
\begin{itemize}
  \item[\checkbox] \texttt{celebration\_bg\_stub.png} or approved equivalent replacement
\end{itemize}

\textbf{Deferred until separately approved}
\begin{itemize}
  \item[\checkbox] full celebration sequence
  \item[\checkbox] confetti / heavy completion VFX
\end{itemize}

\subsection*{H. Technical delivery checks per file}
\begin{itemize}
  \item[\checkbox] native pixel dimensions correct
  \item[\checkbox] no pre-upscale
  \item[\checkbox] transparent background where required
  \item[\checkbox] nearest-neighbor safe pixel structure
  \item[\checkbox] no anti-aliasing
  \item[\checkbox] no blur
  \item[\checkbox] no gradient shading
  \item[\checkbox] naming matches contract
  \item[\checkbox] anchor assumptions preserved
\end{itemize}

\subsection*{I. Source package checks}
\begin{itemize}
  \item[\checkbox] source file included (\texttt{.aseprite} or layered \texttt{.psd})
  \item[\checkbox] export PNGs included
  \item[\checkbox] version labeled clearly
  \item[\checkbox] revision notes included for changed assets
\end{itemize}

\subsection*{J. Explicitly not requested in this handoff}
\begin{itemize}
  \item[\checkbox] alternate themes
  \item[\checkbox] modder assets
  \item[\checkbox] RTL assets
  \item[\checkbox] accessories system
  \item[\checkbox] community pack art
  \item[\checkbox] speculative book covers beyond approved slice
  \item[\checkbox] final atlas packing / optimization
  \item[\checkbox] anything not named in this checklist
\end{itemize}

\subsection*{K. Handoff completion definition}
The artist handoff is complete only when:
\begin{itemize}
  \item[\checkbox] all required approved assets are delivered
  \item[\checkbox] files pass technical import without resizing or rework
  \item[\checkbox] gameplay readability remains intact
  \item[\checkbox] the vertical slice looks cohesive, warm, and production-directed
  \item[\checkbox] replacing placeholders requires texture assignment, not scene redesign
\end{itemize}

\end{document}
